% =====================================================================
% Judul, penulis, dan Abstrak
% Paper semifinal Datathon 2026 · format NeurIPS 2015
%
% Judul dipertahankan dari concept paper penyisihan (konsistensi wajib).
% =====================================================================

\title{Deteksi Dini Kebakaran Hutan dan Lahan Gambut melalui Fusi Multimodal
RGB--Termal dengan Ketahanan Modalitas Hilang Berbasis Drone}

\author{
Argama Vanesa N. Sijabat \\
Matematika, FMIPA \\
Universitas Indonesia \\
\texttt{argamasijabat@gmail.com} \\
\And
Muhamad Erik Setiawan \\
Matematika, FMIPA \\
Universitas Indonesia \\
\texttt{muhamaderiksetiawan4@gmail.com} \\
\AND
Subhan Irsyaduddien Alhaq \\
Matematika, FMIPA \\
Universitas Indonesia \\
\texttt{subhanirsyaduddien@gmail.com} \\
}

\maketitle

\begin{abstract}
Kebakaran hutan dan lahan (karhutla) gambut di Indonesia berulang setiap tahun
dan sulit dideteksi dini karena sistem \emph{hotspot} satelit yang ada tidak
mampu mengonfirmasi titik api kecil yang tersamar asap tebal atau kanopi gambut.
Penelitian ini membangun sistem deteksi dini berjenjang berbasis drone dengan
kontribusi inti berupa model fusi \emph{RGB--termal} berbasis
\emph{cross-attention} yang dilatih dengan \emph{modality dropout}, diperluas
menjadi \emph{reliability gating} adaptif yang disupervisi-diri oleh degradasi
tersintesis. Ketahanan diukur melalui protokol kurva degradasi kontinu pada enam
level, bukan uji biner, dengan uji McNemar berpasangan dan selang kepercayaan
pada setiap angka utama. Fusi terbukti lebih tahan daripada jalur RGB-saja
secara konsisten di tiga partisi data dan unggul atas jalur termal-saja secara
signifikan pada dua di antaranya, tanpa satu pun kejadian api yang terlewat pada
seluruh level degradasi. Kami juga melaporkan hasil yang tidak sesuai hipotesis:
pada kondisi bersih fusi tidak mengungguli jalur RGB-saja, keduanya
menghasilkan prediksi identik pada setiap sampel, \emph{gating} tidak
meningkatkan akurasi, kuantisasi INT8 memperlambat inferensi, dan peringkat
ketahanan berbalik pada partisi uji. Audit label yang kami lakukan sendiri
menemukan satu kelas yang tidak pernah terpelajari, dan dampaknya dilaporkan.
Seluruh komponen terintegrasi menjadi konsol operator terpasang dengan hotspot
satelit langsung. Seluruh citra berasal dari hutan pinus beriklim sedang,
sehingga klaim untuk lahan gambut tropis tidak dibuat.
\end{abstract}
